\def\uadectitulo{Comprensión lectora}
\documentclass[12pt,letterpaper,oneside]{article}
\usepackage[T1]{fontenc}
\usepackage[utf8]{inputenc}
\usepackage[spanish,es-nodecimaldot]{babel}
\usepackage{lmodern}
\usepackage[margin=2.5cm]{geometry}
\usepackage{amsmath,amssymb,graphicx,booktabs,longtable,array,enumitem}
\usepackage[most]{tcolorbox}
\usepackage{tikz}
\usetikzlibrary{babel,arrows.meta,positioning}
\usepackage{setspace}
\usepackage[hidelinks]{hyperref}
\onehalfspacing
\setlength{\emergencystretch}{3em}
\providecommand{\uadectitulo}{Material de estudio}
\providecommand{\uadecautor}{Material de estudio AulaTeX}
\newcommand{\portadaUAdeC}{%
  \hypersetup{pageanchor=false}
  \begin{titlepage}
  \centering
  {\Large Universidad Autónoma de Coahuila\par}
  \vspace{0.7cm}
  {\large Preparación para ingreso a Médico Cirujano\par}
  \vspace{1.8cm}
  {\LARGE\bfseries\uadectitulo\par}
  \vspace{1cm}
  {\large Actividad 1: síntesis y representación conceptual\par}
  {\large Actividad 2: práctica y soluciones justificadas\par}
  \vfill
  {\uadecautor\par}
  {Revisión: 10 de octubre de 2026\par}
  \vspace{0.8cm}
  \begin{minipage}{0.87\linewidth}
  \small Material formativo local. No constituye una guía, plantilla o examen oficial de UAdeC o CENEVAL. La cobertura y dificultad no están certificadas frente a una convocatoria de ingreso. No sustituye orientación ni atención clínica.
  \end{minipage}
  \end{titlepage}
  \hypersetup{pageanchor=true}
  \tableofcontents\clearpage
  \section*{Alcance y procedencia}
  Las actividades siguen las consignas locales conservadas en esta materia. Los reactivos son de práctica y las respuestas son propuestas razonadas, no resultados de un intento del aspirante. Las fuentes de apoyo se identifican al final. Los PDF bibliográficos locales se inventarían por archivo y hash; su disponibilidad no acredita lectura íntegra ni permiso de redistribución.
}
\begin{document}
\portadaUAdeC
\section{Actividad 1: síntesis de investigación}
\subsection{Lectura literal e inferencial}
La lectura literal recupera información explícita. La inferencia relaciona pistas del texto sin añadir hechos externos. La idea principal integra el argumento; un ejemplo es evidencia subordinada, no necesariamente la tesis.
\subsection{Lectura crítica}
Identificar autoría, propósito, tipo de evidencia y límites permite evaluar una afirmación. Una relación temporal o una asociación no prueban causalidad. Es importante separar lo que el texto afirma de lo que el lector considera plausible.
\subsection{Texto original de práctica}
Una biblioteca universitaria amplió su horario durante seis semanas. En ese periodo aumentaron los préstamos y se aplicó una encuesta voluntaria a quienes acudieron por la noche. La mayoría de quienes respondieron valoró el horario extendido. El equipo propuso mantenerlo durante los exámenes, pero aclaró que no podía atribuir el aumento de préstamos solo al horario: también hubo una campaña de difusión. Además, la encuesta no representaba necesariamente a todos los estudiantes. Antes de decidir, recomendó comparar costos y asistencia en otros periodos.
\subsection{Representación conceptual y cuadro de síntesis}
\begin{longtable}{p{0.27\linewidth}p{0.65\linewidth}}
\toprule Concepto & Relación o función\\\midrule
Literal & Qué datos aparecen\\
Inferencia & Qué se deduce de las pistas\\
Evaluación & Qué límites reconoce el autor\\
\bottomrule
\end{longtable}
\begin{figure}[htbp]
\centering
\begin{tikzpicture}[node distance=5mm, every node/.style={draw,rounded corners=2pt,text width=10cm,align=center,font=\small}]
\node (step0) {Literal: Qué datos aparecen};
\node[below=of step0] (step1) {Inferencia: Qué se deduce de las pistas};
\draw[-{Stealth}] (step0) -- (step1);
\node[below=of step1] (step2) {Evaluación: Qué límites reconoce el autor};
\draw[-{Stealth}] (step1) -- (step2);
\end{tikzpicture}
\caption{Esquema didáctico original; las flechas indican orden de lectura o relación conceptual, no causalidad universal.}
\end{figure}
\clearpage\section{Actividad 2: cuestionario original resuelto}
\begin{enumerate}
\item\textbf{¿Durante cuánto tiempo se amplió el horario?}
\begin{enumerate}[label=\Alph*)]
\item Seis meses
\item Seis semanas
\item Un día
\item Un año
\end{enumerate}
\textbf{Respuesta: B.} El texto indica explícitamente seis semanas; las otras duraciones no aparecen.
\item\textbf{¿Cuál es la idea central?}
\begin{enumerate}[label=\Alph*)]
\item Los préstamos nunca dependen de campañas
\item La encuesta representa a todos
\item Se evalúa una ampliación de horario reconociendo límites
\item Los costos ya están calculados
\end{enumerate}
\textbf{Respuesta: C.} El argumento combina resultados preliminares y precauciones, no una conclusión causal ni una decisión final.
\item\textbf{¿Quiénes podían responder la encuesta descrita?}
\begin{enumerate}[label=\Alph*)]
\item Quienes acudieron por la noche y participaron voluntariamente
\item Todos por obligación
\item Solo docentes
\item Solo personal de costos
\end{enumerate}
\textbf{Respuesta: A.} La muestra fue voluntaria entre asistentes nocturnos; no se afirma participación universal.
\item\textbf{¿Por qué no se atribuye el aumento solo al horario?}
\begin{enumerate}[label=\Alph*)]
\item No hubo préstamos
\item Se cerró la biblioteca
\item Todos los datos son falsos
\item También hubo una campaña de difusión
\end{enumerate}
\textbf{Respuesta: D.} La campaña es otra explicación posible y evita atribución causal exclusiva.
\item\textbf{¿Qué decisión propone el equipo?}
\begin{enumerate}[label=\Alph*)]
\item Cerrar definitivamente
\item Mantener el horario durante exámenes
\item Eliminar encuestas
\item Aumentar cuotas de inmediato
\end{enumerate}
\textbf{Respuesta: B.} La propuesta es específica para exámenes; no se anuncia cierre ni cuotas.
\item\textbf{¿Qué limita la generalización de la encuesta?}
\begin{enumerate}[label=\Alph*)]
\item Usar palabras claras
\item Tener seis semanas
\item La selección voluntaria de asistentes nocturnos
\item Contar préstamos
\end{enumerate}
\textbf{Respuesta: C.} Quienes asisten y responden pueden diferir del resto de estudiantes; las otras opciones no explican selección.
\item\textbf{En el texto, 'atribuir' significa:}
\begin{enumerate}[label=\Alph*)]
\item Asignar una causa o explicación
\item Contar libros
\item Pagar costos
\item Cambiar una fecha
\end{enumerate}
\textbf{Respuesta: A.} El verbo vincula el aumento con una posible causa, no con tareas administrativas distintas.
\item\textbf{¿Qué falta comparar antes de decidir?}
\begin{enumerate}[label=\Alph*)]
\item Solo nombres
\item Solo opiniones externas
\item Calificaciones médicas
\item Costos y asistencia en otros periodos
\end{enumerate}
\textbf{Respuesta: D.} La última oración identifica esas comparaciones.
\item\textbf{¿Qué inferencia está respaldada?}
\begin{enumerate}[label=\Alph*)]
\item Todos prefieren la noche
\item La campaña pudo contribuir al aumento
\item El horario no tuvo efecto alguno
\item Los costos fueron nulos
\end{enumerate}
\textbf{Respuesta: B.} Se reconoce una contribución posible, no resultados definitivos ni absolutos.
\item\textbf{¿Cuál afirmación excede la evidencia del texto?}
\begin{enumerate}[label=\Alph*)]
\item Hubo una encuesta voluntaria
\item Hubo difusión
\item La ampliación causó por sí sola todo el aumento
\item Se recomienda comparar periodos
\end{enumerate}
\textbf{Respuesta: C.} El propio texto rechaza esa atribución exclusiva; las restantes afirmaciones son explícitas.
\end{enumerate}

\section{Conclusiones}
El material permite repasar conceptos y contrastar respuestas. La clave no acredita desempeño del aspirante ni aprobación institucional. Deben revisarse la convocatoria y el temario aplicables antes de usarlo como plan definitivo de ingreso.
\clearpage\section{Fuentes de apoyo y alcance de consulta}
Real Academia Española y ASALE. Diccionario panhispánico de dudas, 2.ª edición. Consulta: 10 de octubre de 2026. \url{https://www.rae.es/dpd}\par Consulta de referencia lingüística general; el texto y los reactivos siguientes son originales, no un instrumento estandarizado.\par\medskip
Los originales bibliográficos disponibles se identifican por nombre de archivo y SHA256 en el inventario local. No se atribuyen ediciones, páginas ni lecturas no comprobadas.
\end{document}
